\chapter{المقدمة }
\label{ch:introduction}

تعتبر المناورات الروبوتية، وعلى وجه الخصوص الروبوتات ذات المفاصل الأفقية (\LR{SCARA})، أنظمة ديناميكية شديدة اللاخطية ومتعددة المداخل والمخارج (\LR{MIMO}). إن تحقيق أداء حركي عالي الدقة في هذه الأنظمة لا يقتصر فقط على التصميم الميكانيكي، بل يعتمد بشكل جوهري على عمق الفهم الرياضي للتداخلات الحركية (\LR{Dynamic Coupling})، وقوى القصور الذاتي، وتأثيرات كوريوليس والجاذبية. ومع تزايد متطلبات التطبيقات الحديثة للتعامل مع مسارات معقدة وحمولات متغيرة ($m_L$)، تصبح الحاجة إلى استراتيجيات تحكم متقدمة قادرة على ترويض هذه الديناميكا أمراً حتمياً لا يمكن تجاوزه.

تقليدياً، يتم اللجوء إلى تقنيات التحكم الخطي المستقل لكل مفصل، مثل المتحكمات التناسبية-التكاملية-التفاضلية (\LR{PID}). ورغم بساطتها، إلا أن أداءها يتدهور بشدة عند السرعات العالية بسبب التجاهل التام للتداخلات الديناميكية بين المفاصل. وفي المقابل، تقدم خوارزميات التحكم المعتمدة على النموذج، مثل التحكم بالعزم المحسوب (\LR{CTC})، حلاً رياضياً مثالياً من خلال مبدأ الاستخطاط بالتغذية العكسية (\LR{Feedback Linearization}) لإلغاء اللاخطية تماماً. لكن هذا الإلغاء الجبري يتطلب تطابقاً تاماً بين النموذج الرياضي المحسوب والنظام الفيزيائي الحقيقي، وهو افتراض يصعب تحقيقه في الواقع العملي في ظل وجود شكوك معلمية (\LR{Parametric Uncertainties}) كاحتمالية تغير كتلة الحمل المرفوع، أو وجود قوى احتكاك غير منمذجة بدقة.

لمواجهة تحديات عدم اليقين في النمذجة الرياضية، برزت نظريات التحكم المتين والتكيفي كأقوى الحلول الهندسية في نظرية التحكم الحديثة. يوفر التحكم بالانزلاق (\LR{SMC}) مناعة عالية ضد الاضطرابات، خاصة عند تعزيزه بقوانين وصول متقدمة (\LR{Proportional Rate Reaching Law}) تخفف من ظاهرة الاهتزاز (\LR{Chattering}) وتضمن وصولاً أسياً سريعاً لسطح الانزلاق. من جهة أخرى، يمثل التحكم التكيفي المعتمد على مصفوفة الانحدار (\LR{Regressor-Based ACTC}) ذروة التطور في هذا المجال؛ حيث يمنح الروبوت قدرة ذاتية على "التعلم" وتحديث معاملاته الديناميكية المجهولة لحظياً لتعويض أي تغير في الحمل. ولضمان عدم خروج هذه المعاملات عن السيطرة بسبب ضجيج الحساسات، يتم تحصين خوارزميات التحديث رياضياً باستخدام تقنيات متقدمة مثل عامل التسريب (\LR{$\sigma$-modification}).

يندرج هذا المشروع في صميم هذا التحدي الهندسي؛ حيث لا يكتفي بتطبيق خوارزميات تحكم كصناديق مغلقة، بل يغوص في البنية الرياضية العميقة للروبوت. من خلال تأسيس نموذج ديناميكي دقيق باستخدام صياغة لاغرانج-أويلر، وبناء مصفوفات العطالة والانحدار وتضمين كتلة الحمل بداخلها، ينتقل المشروع تدريجياً لترويض هذه المعادلات برمجياً وتطبيق استراتيجيات تحكم لاخطية وتكيفية متطورة. إن الهدف الجوهري من هذا البحث هو إثبات القدرة على تصميم نظام تحكم متكامل يضمن استقرار ليابونوف (\LR{Lyapunov Stability})، ويحقق تتبعاً مسارياً دقيقاً، ويتكيف بمرونة مع أعقد المتغيرات الفيزيائية، مما يشكل جسراً حقيقياً بين النظرية الرياضية البحتة والتطبيق الصناعي الموثوق.